\section{Do forecasts change for the right reason?}
\label{sec:evidence}

\begin{figure}[!t]
\centering
\includegraphics[width=\linewidth]{exp1_protocol_four_panels.pdf}
\caption{\textbf{When should a forecast change?} An agent forecasts an
unresolved question, then receives one fictional update with search off. Some
updates support the outcome, some oppose it, and some concern the same topic but
do not bear on the outcome. The labels remain hidden; we ask whether the model
moves when the news matters and whether it moves the right way.}
\label{fig:exp1-pipeline}
\end{figure}

Forecasting is not only about assigning a probability; it is also about knowing
which evidence should change it. A model can still move for the wrong reason: a
headline may share a market's topic without changing the event's likelihood. We
therefore begin by asking whether models distinguish evidence that bears on an
outcome from topical resemblance and move in the direction it warrants. Exp.~1
tests this on unresolved markets. Exp.~2 raises the bar: when populations respond
differently, can context select the relationship that governs a new case, and can
local evidence revise it?

\subsection{Exp.~1: can the model separate evidence from topical chatter?}
\label{sec:exp1}

\noindent\textbf{Design.}
We test this before events resolve using 100 unresolved binary
Polymarket questions. Each model first forecasts with web search, then sees one
fictional update with search turned off. Some updates support YES, some support
NO, and others mention the same topic without bearing on the outcome. The model
does not see these labels. Its task is to decide whether to move and, if so, in
which direction. A word-only classifier evaluated on separate markets shows what
surface cues can accomplish; the full sampling design is in
App.~\ref{app:exp1}. Figure~\ref{fig:exp1-pipeline} shows the sequence.
Across nine systems, this yields 30,094 valid forecast updates.

\noindent\textbf{Finding: models first decide whether an update is warranted.}
Figure~\ref{fig:exp1-updating} shows that every system leaves its prior
unchanged more often for topical controls. When models update, they move
$1.9$--$7.5\times$ farther for directional packets across all nine systems. They
also move in the implied direction: directional consistency spans $.887$--$.990$,
compared with $.688$ for the word-only classifier. Human validation and
uncertainty estimates are in App.~\ref{app:exp1-threshold-robustness}.

\begin{figure}[!b]
\centering
\begin{subfigure}[c]{0.525\textwidth}
\centering
\includegraphics[width=.92\linewidth]{exp1_direction_selective_updating.pdf}
\caption{}
\label{fig:exp1-bars}
\end{subfigure}\hspace{0.01\textwidth}
\begin{subfigure}[c]{0.455\textwidth}
\centering
\footnotesize
\setlength{\tabcolsep}{3pt}
\begin{tabular*}{\linewidth}{@{}l@{\extracolsep{\fill}}rccc@{}}
\toprule
\textbf{Model} & \textbf{Size} & \textbf{Abst.} & \textbf{Sens.}$\uparrow$ & \textbf{EHC$_{\geq3}$}$\uparrow$ \\
\midrule
\multicolumn{5}{@{}l}{\emph{Hosted API (parameters undisclosed)}}\\
\micon{claude.png}~Opus~4.8 & --- & 71.4\% & \cellcolor{heatgreen}$7.0\times$ & \cellcolor{heatgreen}.990 \\
\micon{openai.png}~GPT-5.4 & --- & 8.4\% & \cellcolor{heatgreen}$7.5\times$ & \cellcolor{heatgreen}.987 \\
\micon{deepseek.png}~V4-Pro & --- & 44.4\% & \cellcolor{heatyellow}$4.9\times$ & \cellcolor{heatyellow}.977 \\
\midrule
\multicolumn{5}{@{}l}{\emph{Open-weight, descending parameters} $\downarrow$}\\
\micon{qwen.pdf}~2.5-72B & 72B & 24.2\% & \cellcolor{heatorange}$3.2\times$ & \cellcolor{heatyellow}.963 \\
\micon{llamma.pdf}~3.3-70B & 70B & 10.7\% & \cellcolor{heatyellow}$5.7\times$ & \cellcolor{heatyellow}.965 \\
\micon{llamma.pdf}~3.1-70B & 70B & 2.5\% & \cellcolor{heatyellow}$4.5\times$ & \cellcolor{heatyellow}.961 \\
\micon{qwen.pdf}~2.5-32B & 32B & 11.1\% & \cellcolor{heatorange}$2.8\times$ & \cellcolor{heatyellow}.964 \\
\micon{qwen.pdf}~2.5-14B & 14B & 26.1\% & \cellcolor{heatred}$2.2\times$ & \cellcolor{heatorange}.943 \\
\micon{llamma.pdf}~3.1-8B & 8B & 1.1\% & \cellcolor{heatred}$1.9\times$ & \cellcolor{heatred}.887 \\
\bottomrule
\end{tabular*}
\caption{}
\label{fig:exp1-table}
\label{tab:consistency}
\end{subfigure}
\caption{\textbf{Models move much more for evidence than for topical
chatter, and usually move in the implied direction.} \textbf{(a)} Average
revision after supporting, contrary, and topical updates. \textbf{(b)} The
table summarizes how often models leave topical updates alone, how much farther
they move for directional evidence, and how often a substantial move has the
right sign. Intervals and robustness checks are in
App.~\ref{app:exp1-threshold-robustness}.}
\label{fig:exp1-updating}
\end{figure}

\begin{figure}[!t]
\centering
\includegraphics[
  width=.90\linewidth,
  trim=11pt 14.5pt 3.5pt 7pt,
  clip
]{exp2_coin_city_prompt_arms.pdf}
\caption{\textbf{Coin City turns contextual reasoning into a paired
comparison.} Every arm shares the query and numerical cases. The context arm
adds a description connecting City C to one displayed media environment; the
opposite-context arm changes only that sentence. Target cases then accumulate,
allowing local evidence to confirm or correct the contextual choice.}
\label{fig:coin-city-prompt}
\end{figure}

\noindent\textbf{\emph{What Exp.~1 establishes.}}
\emph{Across the tested systems, forecasts distinguish outcome-relevant evidence
from topical chatter and move in the direction implied by the evidence.}
This does not yet show whether the model knows whose response pattern makes that
evidence predictive. Coin City exposes that hidden choice.

\subsection{Exp.~2: let context select a relationship, then challenge it}
\label{sec:exp2}

\noindent\textbf{Coin City makes the process ambiguity observable.}
Figure~\ref{fig:world-knowledge-evidence} began with a familiar ambiguity: the
same two heads can imply different forecasts depending on what process the
forecaster believes generated them. Coin City turns that intuition into a
controlled test. The same campaign story moves polls strongly in City A but
weakly in City B; context identifies the better analogue for target City C. If
the model uses that link, changing only the context should change its forecast,
while direct City C cases should eventually override a misleading cue. Each
episode supplies noisy A/B examples, a one-sentence link, and zero to four target
cases.
Each city's poll response follows
\[
  \Delta p_{ij}=g_j x_{ij}+\epsilon_{ij},
  \qquad \epsilon_{ij}\sim\mathcal{N}(0,5^2),
  \label{eq:coin-city-dgp}
\]
where \(x_{ij}\) is signed campaign news and \(g_j\) is that city's stable
sensitivity. The prompt never states \(g_j\); the model has to infer it from
examples.
Figure~\ref{fig:coin-city-prompt} shows the shared query, reference cities, and
target-city context in one illustrated arm.

We compare five versions of the same problem: target cases alone; references
without a link; the correct link; the opposite link; and episode-varying KIV/ZOR
labels whose meaning must be learned from the references. Forecast error tells us
whether context helps. The model's implied response to campaign news tells us
which city it is treating as the analogue. App.~\ref{app:coin-city-results}
gives the full model roster and sampling design.

\noindent\textbf{Finding: the forecast follows the selected analogue.}
Figure~\ref{fig:coin-city-results} shows that correct context lowers error before
any City C cases for every system; changing only that sentence reverses the selected
regime. The implied response correlates $.51$--$.86$ with the named reference
after controlling for the true regime, while the other reference is near zero.
As City C cases accumulate, both the contextual advantage and inversion harm
shrink: forecasts shift from context toward diagnostic local behavior.

\begin{figure}[!t]
\centering
\includegraphics[width=\linewidth]{exp2_coin_city_results.pdf}
\vspace{3pt}
\begin{minipage}{\linewidth}
\centering
\footnotesize
\setlength{\tabcolsep}{4pt}
\resizebox{\linewidth}{!}{%
\begin{tabular}{l r rrrr rr}
\toprule
& & \multicolumn{4}{c}{Forecast MAE [95\% CI]} & \multicolumn{2}{c}{$\rho(\hat g,g_C)$} \\
\cmidrule(lr){3-6}\cmidrule(lr){7-8}
Model & Paired $n$ & Target only & No C context & $+$ C context &
$+$ wrong C context & No context & $+$ context \\
\midrule
\micon{claude.png}~Opus~4.8 & 250 & 2.32 [2.20,\,2.46] & 2.65 [2.41,\,2.89] & \textbf{1.82} [1.64,\,2.02] & 4.20 [3.87,\,4.54] & $-.06$ & \textbf{.70} \\
\micon{claude.png}~Opus~5 & 250 & 2.36 [2.12,\,2.59] & 2.62 [2.38,\,2.84] & \textbf{1.76} [1.57,\,1.99] & 4.25 [3.92,\,4.57] & $-.01$ & \textbf{.62} \\
\micon{openai.png}~GPT-5.6 & 250 & 3.60 [3.28,\,3.94] & 2.52 [2.33,\,2.75] & \textbf{1.60} [1.42,\,1.77] & 4.23 [3.90,\,4.55] & $-.04$ & \textbf{.68} \\
\micon{deepseek.png}~V4-Pro & 250 & \textbf{2.45} [2.27,\,2.61] & 2.84 [2.55,\,3.15] & 2.47 [2.19,\,2.77] & 4.36 [4.01,\,4.74] & $-.03$ & \textbf{.52} \\
\micon{kimi.png}~K3 & 250 & 3.35 [3.04,\,3.67] & 2.55 [2.32,\,2.78] & \textbf{1.89} [1.66,\,2.13] & 4.31 [3.98,\,4.62] & $-.02$ & \textbf{.59} \\
\micon{gemini.png}~3.6 Flash & 250 & 2.53 [2.35,\,2.73] & 2.55 [2.31,\,2.80] & \textbf{2.17} [1.89,\,2.48] & 4.29 [3.94,\,4.66] & $.08$ & \textbf{.56} \\
\bottomrule
\end{tabular}
}
\end{minipage}
\caption{\textbf{Context improves forecasts before target cases and helps assign
the response pattern across six systems.} Giving the correct analogue lowers
error before City C has a history; naming the wrong analogue raises it. As direct
City C cases accumulate, forecasts rely less on either contextual cue. Full
results by model are in App.~\ref{app:coin-city-results}.}
\label{fig:coin-city-results}
\end{figure}

\subsubsection{Infer a convention with no stable semantics}
\label{sec:novel-mapping}

The context result still admits a weaker explanation: familiar descriptions such
as ``national news'' might cue a response without the model inferring the
relationship. KIV and ZOR remove that route. Their meanings change in every
episode, so the model can use them only by inferring from the reference cases
which label names which response pattern. Twelve of fifteen systems recover the
convention before any target cases. When the two populations are made harder to
distinguish, performance falls but the ordering of models remains similar, as an
induction account predicts (App.~\ref{app:coin-city-robustness}).

\begin{wrapfigure}[23]{r}{0.43\textwidth}
\vspace{-1.8em}
\centering
\includegraphics[width=\linewidth]{exp2_symbol_decoding.pdf}
\vspace{-1.0em}
\caption{\textbf{Larger models infer an episode-local convention, then yield
to direct evidence.} On 88 new episodes, four larger Qwen and Llama checkpoints
(14--72B) recover the changing KIV/ZOR mapping; three smaller checkpoints
(4--8B) struggle.}
\label{fig:coin-city-symbol-decoding}
\vspace{-1.8em}
\end{wrapfigure}

This behavior shows that models can use a newly induced convention, but not
whether they represent that convention internally or whether it helps produce
the forecast.

Figure~\ref{fig:coin-city-symbol-decoding} asks the first question. Across seven
open checkpoints, the four larger models carry the episode's inferred KIV/ZOR
meaning in their internal state, while the three smaller models struggle. That
state follows the contextual relationship before City C has a history, then
yields as direct City C evidence identifies the local response. The representation
therefore tracks which source of information should govern the forecast, rather
than a fixed meaning of either label.

Decoding is still correlational. Figure~\ref{fig:coin-city-causal-patch} asks the
causal question by swapping the label-state representation between otherwise
matched Qwen3-14B prompts. Before City C provides its own history, the intervention
changes both the analogue the model names and its numerical forecast. After four
City C cases, the same swap has no reliable effect. Across episodes, the shift is
\(+.333\) before target cases and \(-.035\) after four, the latter inside the
random-swap range. This timing is the point: the induced convention helps produce
the forecast when context is needed, then gives way when direct evidence reveals
the relationship (App.~\ref{app:coin-city-causal-patch}).

\noindent\textbf{\emph{What Exp.~2 establishes.}}
\emph{Models use context to select which population's response governs a
forecast and revise that choice when target evidence arrives. Twelve of fifteen
systems infer a convention with no stable semantics; larger open models encode
that episode-local mapping, and changing it causally changes Qwen3-14B's forecast
when local evidence is absent.}

\begin{figure}[!t]
\vspace{-0.4em}
\centering
\includegraphics[width=\linewidth]{exp2_causal_patch.pdf}
\vspace{-0.7em}
\caption{\textbf{Changing the induced convention changes the forecast before
target evidence arrives, but not after.} In the illustrated pair, swapping the
internal representation moves the forecast from 51.3 to 54.0 and changes which
reference city the model names; reversing the swap reverses the movement. Once
four target cases reveal City C's own behavior, arbitrary swaps are no more
influential than the intervention. Full controls and across-episode results are
in Table~\ref{tab:coin-city-causal-patch}.}
\label{fig:coin-city-causal-patch}
\vspace{-1.0em}
\end{figure}

Section~4 asks whether training carries this form of relationship use into a new
world.
